\documentclass[12pt,a4paper,twoside]{ipb}

% comentar caso seja uma disseração de mestrado
%\usepackage{projei}

\usepackage[english]{babel}
\graphicspath{{./images/}}
\usepackage{listings} % incluir listagens
\usepackage{url} % typeset URL's
\usepackage[colorlinks=true,
                    urlcolor=black, %blue
                    linkcolor=black,
                    citecolor=black, %cor das citações
                    bookmarks=true,
                    pdfstartview=FitH]{hyperref}

% Pode ser usado biblatex
\usepackage[style=ieee,backend=biber]{biblatex}
\addbibresource{lib/refs.bib}
\usepackage{todonotes}
\usepackage{lipsum}
\usepackage{tabulary}
\usepackage{algorithm}
\usepackage{boxedminipage}
\usepackage{algpseudocode}
\usepackage{caption}
\usepackage{booktabs}
\usepackage{array}
\usepackage{color,soul}
\usepackage{pdfpages}

%%%%%%%%%%%%%%%%%%%%%%%%%

\title{XXXX}

\author{name student}
%\authnum{1}
%\secondauthor{Student name - Student number}
%\secauthnum{2}

\supervisor{Prof. XXXX}
\cosupervisor{Prof. XXXX}

% Para definir o ano letivo
\courseyear{202X-202X}


%\cosupervisor{Nome do co-orientador}
%\firstreader{Nome do arguente}
% Para definir o ano letivo
%\courseyear{1234}
% Para nao mostrar a lista de tabelas
\tablespagefalse

\makeglossaries
\loadglsentries{acronym}

\begin{document}
	
	
	% Coloca a capa, primeira pagina e outros
	
	\beforepreface
	
	%\cleardoublepage
	
	\prefacesection*{Dedication}
	\input{chapters/dedication}
	
	%\cleardoublepage
	
	\prefacesection*{Acknowledgment}
	\input{chapters/acknowledgment}
	
	%\cleardoublepage
	
	\prefacesection*{Abstract}
	\input{chapters/abstract}

	\prefacesection*{Resumo}
	\input{chapters/resumo}

	%\cleardoublepage

	
	% Coloca indices
	\afterpreface
	%\cleardoublepage
	\printglossary[type=\acronymtype,nonumberlist,title={Acronyms}]
	\glsaddall
	
	%\listoftodos
	
	\bodystart
	
	
	% Capitulos do documento
	\cleardoublepage
	\input{chapters/cap01}
	\cleardoublepage
	\input{chapters/cap02}
	\cleardoublepage
    % add more if necessary
    
	
	%% estilo de referências. outros valores posíveis são 'plain' e 'abbrv' apalike
	%\bibliographystyle{plain}
	%% listagem de referências
	%\bibliography{lib/refs}
	
	%  Caso seja usado biblatex
	\printbibliography
	
	
	% Apêndices
	\appendix
	
	%http://tex.stackexchange.com/questions/59572/custom-page-numbering-for-appendix
	\pretocmd{\chapter}{%
		\clearpage
		\pagenumbering{arabic}%
		\renewcommand*{\thepage}{\thechapter\arabic{page}}%
	}{}{}
	
	\input{chapters/apendice01}
	
\end{document}